\documentclass[../main.tex]{subfiles}
\begin{document}
%==========================================TIÊU ĐỀ TẬP========================================
\begin{table}[h]
  \begin{adjustwidth}{-1cm}{}
    \begin{tabular}{>{\centering\arraybackslash}p{6cm}|>{\raggedright\arraybackslash}p{12.5cm}}
      \multirow{5}{6cm}
      % Cột bên trái: ÉPISODE và số tập
      {\\[-40pt]
        \flaregothic\fontsize{20pt}{20pt}\selectfont \centering
        \color{eptype}ÉPISODE\color{black}\\
        \burbank\fontsize{72pt}{72pt}\selectfont
        \color{epnum}66\color{black}\\[6pt]
        \cafeteria\fontsize{20pt}{20pt}\selectfont\color{epnum}(5-10)\color{black}
        \\[18pt]
        % \flaregothic\fontsize{14pt}{14pt}\selectfont
        % \color{epattr}BIENVENUE, \uline{\color{Banpire} Mel} !
      }
       &
      % Dòng thứ nhất: tên tập tiếng Nhật
      {\fotskip\fontsize{18pt}{18pt}\selectfont お\ruby{願}{ねが}いがあるの…
      }                                                           \\[6pt]
      % Dòng thứ hai: tên tập tiếng Anh
       & {\cafeteria\fontsize{18pt}{18pt}\selectfont Caltrops}                                                               \\[4pt]
      % Dòng thứ ba: tên tập tiếng Việt
       & {\vnmsans\fontsize{18pt}{18pt}\selectfont Cây đèn thần}                                                             \\[8pt]
      % Dòng thứ tư: Nhân vật xuất hiện
       & {\caviar{\fontsize{14pt}{14pt}\selectfont Nhân vật: \emp{kuon1} \emp{suisei} \emp{matsuri} \emp{haato} \emp{luna}}}
      \\[8pt]
      % Dòng cuối cùng: Ngày phát hành
       & {\continuum\fontsize{16pt}{16pt}\selectfont Ngày phát hành: 09/08/2020}                                             \\[18pt]
    \end{tabular}
  \end{adjustwidth}
\end{table}
%=========================================TRUYỆN CHÍNH========================================
\color{black}

% Hộp tóm tắt
\txtbx[2]{{\color{vol} \uline{Tóm tắt:}} Matsuri, Kuon và Haato tỉnh dậy trong văn phòng đầy chông sắt do Haato bày trò. Họ tìm thấy cây đèn thần, nhưng thay vì thần đèn, Suisei xuất hiện với vai Akinator và ban điều ước theo cách\dots\ sai lệch. Matsuri ước ramen nhưng bị Haato cướp lời thành sư tử, rồi được đưa đến một quán ramen thật. Haato ước dẹp chông nhưng bị Matsuri cướp lời, rồi bị đưa đến Nam Cực. Cuối cùng, Kuon gọi Luna ra, dùng ba điều ước để dọn dẹp văn phòng, về nhà ăn tối cùng gia đình, và\dots\ vô tình để quên Haato đang lạnh cóng ở Nam Cực.}

Ngày đầu tiên của tuần lễ Obon, trời trong xanh, nắng nhẹ, và văn phòng \hll\ vắng tanh, vì hầu hết mọi người đều về quê thăm gia đình. Thế nhưng, trong căn phòng im lìm ấy, Matsuri đang say sưa với một giấc mơ thơm ngon. Cô thấy mình ngồi trước một tô mì ramen bốc khói, sợi mì dai dai, nước dùng đậm đà, thịt cháy cạnh vàng ươm.

Nhưng rồi, bát mì trong mơ đã hết. Và đúng lúc đó, cô tỉnh dậy.

Điều đầu tiên đập vào mắt cô không phải là trần nhà quen thuộc, cũng chẳng phải ánh nắng ban mai, mà là một khung cảnh kỳ quặc đến khó tin. Chi chít những chông sắt sắc nhọn, được rải đều khắp sàn nhà như một bãi mìn thu nhỏ. Cô choàng dậy, hoang mang thốt lên, giọng đầy vẻ không tin nổi: ``Cái quái gì\dots''

\img{5-10a}

Bên cạnh cô, trên một chiếc sofa cũng đang nằm ở giữa phòng, Kuon cũng vừa tỉnh dậy với vẻ mặt bối rối tương tự. Cậu ngồi dậy, dụi mắt và nhìn xung quanh, cố gắng hiểu chuyện gì đang xảy ra.

Chợt, một giọng nói quen thuộc vang lên từ phía cửa, đầy vẻ vui vẻ đến kỳ lạ: ``Hai người tỉnh rồi à?''

Matsuri nhận ra đó là Haato. Cô nàng đang đứng ở cửa, tay cầm một tấm bảng, khuôn mặt rạng rỡ như thể vừa hoàn thành một dự án nghệ thuật. Cô hỏi trong sự hoang mang: ``Haato-chan? Bà làm gì ở đây vậy? Sao lại có chông sắt khắp sàn thế này?''

Kuon cũng thắc mắc, giọng đầy nghi ngờ: ``Chuyện gì đang xảy ra ở đây? Sao Mat-chan lại tọt ra tận đó? Và ai đã rải đống chông này xuống?''

Haato vui vẻ đáp, không một chút áy náy, như thể đang khoe một thành tích đáng tự hào: ``Tôi vừa bê cái sofa này ra đây đấy. Còn cái kia thì không bê được vì Kuon-senpai nặng quá. Tôi đã thử mấy lần nhưng không nhấc nổi.''

Cậu Tổng ngạc nhiên: ``Thế là thế nào? Tự dưng bê cái sofa ra giữa nhà làm gì? Với lại\dots\ đống chông sắt này có ý nghĩa gì?''

Matsuri nhìn xuống sàn, tay chỉ về những chiếc chông sắc nhọn đang lấp lánh dưới ánh đèn: ``Đống dưới sàn là\dots\ chông sắt đấy à? Bà định biến văn phòng thành pháo đài phòng thủ hay gì?''

Haato suy nghĩ một lúc, rồi vui vẻ trả lời, tay cầm tấm bảng - vốn là đạo cụ mà Okayu đã từng dùng trước đó - giơ lên về phía họ: ``Chuẩn luôn! Tại tôi thấy vui quá nên đã rải chúng ra đấy! Không khí Obon mà, phải có chút gì đó đặc biệt chứ!''

\img{5-10b}

Kuon nghiêng đầu, ``hả'' một tiếng, rồi nói trong sự ngán ngẩm, giọng của một người đã quá quen với sự vô lý: ``Em rảnh rỗi quá à? Không có việc gì làm nên bày trò hả?''

Haato ``ưm'' một tiếng một cách vui tươi và nói tiếp với sự khoái chí, hai tay đặt lên má, trông như một đứa trẻ vừa hoàn thành một trò chơi: ``Tiện thể thì giờ em cũng bị kẹt ở đây với hai người luôn rồi. Em đã rải chông ra, nhưng quên mất là mình cũng phải đi qua đó.''

Matsuri và Kuon đồng thanh đáp lại, mặt cau có, giọng như một bản hòa ca bất đắc dĩ: ``Cái đó không phải là thứ mà bà/em nên thấy tự hào đâu!''

Trong lúc Kuon đang tiếp tục nghĩ ra câu mắng cho Haato vì sự ngớ ngẩn của cô, cậu bất chợt khựng lại. Mắt cậu dừng lại ở một vật thể trên bàn giữa phòng. ``Mà này\dots\ có cái gì trên bàn kia?''

Cả ba người nhìn về phía mà cậu chỉ. Trên bàn, nằm giữa đống giấy tờ lộn xộn, là một cây đèn đồng cổ kính, với họa tiết chạm khắc tinh xảo và một cái vòi cong cong, trông giống hệt cây đèn thần trong chuyện Aladdin.

Matsuri ngạc nhiên, bước lại gần, tay đưa ra: ``Đây chẳng phải là cái đèn ban điều ước sao? Nhìn giống hệt trong phim ấy!''

Kuon cũng đồng tình, gật đầu, ánh mắt đầy nghi ngờ: ``Cũng có thể. Nhìn như cái đèn mà bốn mươi tên cướp muốn lấy đây. Mà ai lại để nó ở đây giữa đống chông sắt này?''

\img{5-10c}

Matsuri cầm lấy cây đèn đó và xoa nhẹ vào thân đèn, như trong những câu chuyện cổ tích. Ngay lập tức, một làn khói xanh đặc quánh bốc ra từ vòi đèn, lan tỏa khắp phòng. Từ trong làn khói ấy, một bóng người hiện ra. Hai tay giơ lên đầy khí thế, giọng vang vọng: ``\textbf{Hoshimachi\~{}!!}''

Nhưng khi làn khói tan đi, bóng người đó không phải là một vị thần linh huyền bí, mà là một cô gái quen thuộc với mái tóc xanh và đôi mắt sao chổi rực sáng: Hoshimachi Suisei.

Matsuri thở dài, hơi thất vọng, tay hạ xuống: ``\dots\ Ơ không, chỉ là Sui-chan thôi. Tưởng là thần đèn thật chứ.''

Kuon lên tiếng thắc mắc, giọng đầy nghi ngờ: ``Em ấy làm gì trong này vậy? Có ai nhốt em ấy vào đèn không?''

Haato nhí nhảnh đáp, giọng như một người vừa khám phá ra một sự thật hiển nhiên: ``Có lẽ là để giúp chúng ta. Logic là vậy; ai lại nhốt mình vào đèn làm gì nếu không phải để giúp người khác?''

Suisei ngó quanh một lúc, đôi mắt sáng lên khi thấy ba người họ đang ngồi trên sofa giữa biển chông sắt. Cô vui vẻ giơ một tấm bảng điện tử (?) có dòng chữ ``Cây đèn thần đoán tên người'' cùng với các lựa chọn như ``Có'', ``Không'', ``Không biết'', ``Chắc có'' và ``Chắc không'', và hào hứng nói: ``Cần chị giúp gì không?''

Matsuri nhìn bảng một lúc rồi nhún vai, giọng đầy vẻ thất vọng khi nhận ra trò chơi: ``À không, đây là Akinator\footnote{Trò chơi thần đèn đoán nhân vật (chủ yếu là các YouTuber) dựa theo gợi ý. Các lựa chọn có ở hình tại trang \pageref{akinator}.} thì có.''

\begin{figure}[H]
  \centering
  \includegraphics[width=12.5cm]{../../../../Image/Book5/5-10d.png}
  \caption{Các lựa chọn bao gồm: ``Có'', ``Không'', ``Không biết'', ``Chắc có'' và ``Chắc không'', giống như một trò chơi Akinator thật.}
  \label{akinator}
\end{figure}

Nàng Cổ động sau đó cũng chống nạnh, than thở với giọng đầy bực bội: ``Tôi đang bị kẹt ở đây này, chán thật chứ! Sàn đầy chông, sofa ở giữa, mà lại còn bị nhốt cùng Haato-chan nữa.''

Kuon thở dài ngán ngẩm, đưa tay xoa trán: ``Chỉ vì có kẻ nào đó nghĩ rằng nó vui. Mà kẻ đó đang đứng ngay trước mặt chúng ta đây.''

Cả hai cùng liếc sang Haato, người bày ra cái trò này. Nàng Song tính lúng túng nhìn đi chỗ khác, giả vờ không biết gì, hai tay chắp sau lưng, miệng huýt sáo nhẹ.

``Với cả, bọn anh không gọi em lên để chơi trò đoán nhân vật. Bọn anh cần em giúp,'' Kuon nói, xua cái bảng chọn kia đi với vẻ mặt đầy bất lực. ``Em có thể giúp bọn anh ra khỏi đây không?''

Suisei đặt tay lên ngực, khẽ nghiêng đầu đầy thần bí, như một vị thần đèn thực thụ: ``Vậy thì, em sẽ ban một điều ước cho mọi người!''

Cả Kuon và Haato nghe vậy liền phấn khởi, mắt sáng lên như hai đứa trẻ vừa được hứa tặng quà.

Kuon gật đầu: ``Thế thì tốt quá. Ít nhất cũng có người biết giúp đỡ.''

Haato vỗ tay: ``Tuyệt! Giờ thì ta có thể ra khỏi đây rồi! Cuối cùng cũng thoát khỏi đống chông sắt này!''

Trong khi Haato đang còn phấn khởi, Matsuri bỗng nhiên trở nên trầm ngâm. Khuôn mặt cô, vốn dĩ lúc nào cũng rạng rỡ, giờ đây ủ rũ xuống, mắt nhìn xa xăm, miệng thở dài não nề. Cả Haato và Kuon đều nhận ra sự thay đổi kỳ lạ này.

``Matsuri-chan?'' Haato ngừng lại, nghiêng đầu hỏi với vẻ tò mò. ``Sao thế? Bà đang nghĩ gì mà mặt mày tiu nghỉu thế?''

Kuon cũng tiếp lời, giọng có chút lo lắng, vì Matsuri hiếm khi im lặng lâu như thế này: ``Đói đến mức không nói được à?''

Nàng Cổ động chợt kêu lên đầy đau khổ, hai tay ôm bụng, giọng như một người vừa mất đi thứ quý giá nhất: ``Bây giờ tôi đang rất thèm ramen! Tôi đói quá! Cái bát mì trong mơ vừa nãy cứ ám ảnh tôi mãi!''

Cả Haato và Kuon đồng loạt hét lên, giọng đầy bất lực: ``\textbf{Bà/em có thể chọn cái gì quan trọng hơn được không!? Đang mắc kẹt giữa đống chông sắt mà lại nghĩ đến ramen!}''

Matsuri vừa mở miệng định nhờ nàng đèn thần gọi cho cô một bát mì ramen (ラーメン), thì Haato lập tức bịt miệng cô lại. Cô nàng Song tính vội vã chen ngang bằng giọng gấp gáp, như thể đang cứu một người sắp rơi xuống vực: ``Bà ấy muốn có một con sư tử (ライオン) đấy!''

Chỉ trong chớp mắt, một con sư tử đực khổng lồ xuất hiện giữa phòng. Nó cao lớn, bờm vàng óng, ánh mắt kiêu hãnh, nhưng lại đang đứng giữa một văn phòng đầy chông sắt một cách bình thản, như thể đó là chuyện hiển nhiên. Còn Suisei thì biến mất trong làn khói xanh, để lại một khoảng không trống rỗng.

Matsuri gục đầu xuống, nức nở với giọng đầy thất vọng, như một đứa trẻ bị cướp mất kẹo: ``Ramen của tôi\dots\ Tôi đã có thể gọi một bát ramen ngon lành, mà lại thành sư tử\dots\ Tôi ghét cuộc đời này\dots''

Kuon nhăn mặt, nhìn Haato đầy khó hiểu, giọng như một người đang cố gắng giải thích một hiện tượng vật lý phi lý: ``Tự dưng lại đòi một con sư tử làm gì vậy? Nó có thể ăn thịt chúng ta trước khi chúng ta kịp rời khỏi chỗ này đấy!''

Haato chắp tay trước ngực, cười gượng, mắt nhìn xuống sàn như một đứa trẻ bị bắt quả tang: ``Tại lúc đó em hoảng quá, chả biết nói gì cả! Em chỉ nghĩ đến ramen, nhưng lại sợ rằng nếu nói ramen, cậu ta sẽ biến ramen thành thứ gì đó kỳ quặc, nên em đành nói bừa!''

Cậu Tổng nhướng mày, thở dài bất lực, giọng đầy mệt mỏi: ``Thế thì thà đừng nói còn hơn! Rải chông khắp văn phòng là chưa đủ hay sao?''

Haato lập tức quát lên, giọng như một cơn bão nhỏ: ``\textbf{Thì đã bảo là em xin lỗi mà! Em biết là sai rồi!}''

Trong lúc cô nàng còn đang rối rít xin lỗi Kuon với đôi mắt ngấn nước, con sư tử bất chợt gầm lên một tiếng mạnh mẽ, vang dội, khiến cả căn phòng rung lên và ba người đồng loạt giật mình.

Matsuri ngẩng đầu nhìn nó, mắt vẫn còn đỏ hoe, rồi đột nhiên cất giọng hỏi với vẻ đầy thất vọng, như thể cô không có hy vọng gì với con sư tử kia: ``Sư tử nè, mày có thể cho tao một bát ramen được không? Mì kiểu gì cũng được, đói lắm rồi\dots''

Kuon cau mày, liếc nhìn con sư tử rồi lắc đầu nói, giọng đầy hoài nghi của một người đã thấy quá nhiều điều kỳ lạ: ``Anh không nghĩ là nó biết nói đâu, huống gì--''

Câu nói chưa dứt, con sư tử mà Suisei vừa triệu ra đã đứng dậy. Nó bước những bước nặng nề về phía kệ bếp, cầm lấy một cái bát trên giá, rồi quay lại, cất giọng trầm ấm, đầy nam tính: ``Có ngay!''

Ba người sững sờ nhìn chằm chằm vào cái bát mà con sư tử đang đưa cho Matsuri. Cái bát trống không. Không có lấy một giọt nước dùng, không một sợi mì, chỉ có một chiếc bát trắng tinh, nhẵn nhụi.

\begin{figure}[H]
  \centering
  \begin{subfigure}[t]{0.45\textwidth}
    \centering
    \includegraphics[width=8cm]{../../../../Image/Book5/5-10f1.png}
    \caption{Con sư tử đứng phắt dậy, chưng ra một cái bát.}
  \end{subfigure}
  \quad
  \begin{subfigure}[t]{0.45\textwidth}
    \centering
    \animategraphics[autoplay,loop,width=8cm]
    {30}{../../../../Image/Book5/5-10f2/5-10f2.}{1}{60}
    \caption{Bên trong cái bát đó là Suisei (tí hon) nhảy ra.}
  \end{subfigure}
\end{figure}

Ấy thế mà, vì quá đói và quá tuyệt vọng, Matsuri liền reo lên đầy phấn khích, mắt sáng rực như vừa thấy kho báu: ``Ramen kìa! Cuối cùng cũng có ramen!''

Cô nàng giơ tay định chộp lấy cái bát, nhưng Kuon nhanh chóng đỡ trán, giọng đầy thương hại nhưng cũng không kém phần bất lực: ``Mat-chan, đó chỉ là cái bát trống thôi. Không có mì, không có nước dùng, không có gì hết. Đừng có tự lừa dối bản thân.''

Phải mất vài giây, cô nàng mới dần nhận ra có gì đó không ổn. Matsuri chớp mắt vài lần, cúi xuống nhìn cái bát rỗng, rồi ngẩng lên nhìn con sư tử, rồi lại nhìn cái bát: ``Hả? Từ từ đã\dots\ Đây là cái bát trống mà! Tôi tưởng có ramen chứ!''

Ngay khi cô còn đang hoang mang, một làn khói xanh bất ngờ bốc lên từ cái bát trống. Đó là khói của cây đèn thần. Suisei từ trong đó chui ra, mặt rạng rỡ như chưa có chuyện gì xảy ra, vẫn giữ nguyên nụ cười và tấm bảng trên tay.

Ba người đồng thanh thốt lên, giọng như một bản hòa ca bất đắc dĩ: ``Đó là Sui-chan mà!''

Cô nàng ``trong đèn'' lại tiếp tục hỏi với giọng vui vẻ như trước, và vẫn cái bảng chọn kia xuất hiện trước mặt họ, như thể cô chưa từng biến mất: ``Cần em giúp gì không?''

Matsuri và Kuon nhìn nhau, rồi cùng lẩm bẩm với giọng đầy thất vọng, như hai người vừa nhận ra rằng mình đang bị lừa: ``Lại là Akinator à\dots\ Không có thần đèn nào thật cả, chỉ có trò chơi đoán nhân vật thôi.''

Lần này, Haato nhanh chóng lao lên trước, mắt sáng rực, đầy quyết tâm. Cô đã không ăn gì từ sáng, và cô không muốn lặp lại sai lầm của Matsuri: ``Tốt rồi! Hãy dẹp hết đống chông này giùm em đi!''

Tuy nhiên, ngay khi dứt câu, Matsuri - trong đầu vẫn đang nghĩ đến hàng nghìn kiểu dáng bát mì ramen khác nhau - liền nhảy vào, bắt chước hành động của Haato vừa nãy. Cô bịt miệng Haato lại rồi nói thật nhanh, như một người đang cố gắng tranh giành điều ước cuối cùng: ``\textbf{Mau đưa tôi ra quán ramen ngay!}''

Matsuri vừa dứt lời, một luồng sáng xanh bỗng chốc bao trùm lấy cô. Cô cảm nhận được một lực hút mạnh mẽ, như thể có ai đó đang kéo cô vào một cơn lốc xoáy phát ra từ cây đèn.

``K-Khoan đã!? Tôi chưa kịp chuẩn bị tinh thần--!!''

Cô còn chưa kịp phản ứng, cơ thể đã bị hút vào một cơn lốc xoáy phát ra từ cây đèn. Trong nháy mắt, cô biến mất hoàn toàn khỏi căn phòng đầy chông sắt, để lại Haato và Kuon đứng nhìn nhau trong sự ngơ ngác.

Haato chớp mắt, nhìn vào khoảng không nơi Matsuri vừa đứng, rồi quay sang Kuon: ``Ể? Gì thế này?''

Haato vội vàng quay sang Suisei, người vẫn đang lơ lửng giữa không trung với vẻ mặt thản nhiên như chưa có chuyện gì xảy ra, tay vẫn cầm cái bảng chọn lựa Akinator. Cô nàng Trung học chạy lại gần, tay chỉ vào ``nàng thần đèn'' với giọng đầy phẫn nộ: ``Khoan đã, tại sao lại là cô ấy!? Em là người nói trước mà! Em còn phát âm rõ ràng nữa! Rõ ràng em đã nói `dẹp chông' rất rành mạch! Còn Matsuri-chan thì nhảy vào sau cơ mà!?''

Suisei chỉ nhún vai, giọng đều đều như một câu trả lời tự động: ``Cô ấy đã ước trước. Không thể thay đổi thứ tự ước.''

Kuon nhìn Haato một lúc, rồi chỉ đơn giản nhún vai và giải thích với giọng của một người đã hiểu rõ luật chơi: ``Bởi vì em đã ước một lần rồi. Con sư tử vừa nãy không tự nhiên xuất hiện đâu.''

Haato đơ người, mắt mở to: ``Hả? Em ước gì cơ? Em chỉ nói `dẹp chông' thôi mà!''

Kuon chỉ trỏ về phía con vật to lớn vẫn còn đứng đó, vô tư gặm móng vuốt của mình như thể đang ở trong sở thú - ``Sui-chan, hay Akinator gì đó vừa nãy nói đấy, chỉ ban một điều ước cho mỗi người, nên điều ước của em đã xong ngay từ lúc em bảo `con sư tử' rồi. Nó coi đó là điều ước của em. Dù em không cố ý.''

Haato há hốc miệng với lời giải thích của cậu. Cô đưa tay ôm đầu, vẻ mặt đầy hối tiếc, giọng như một người vừa phát hiện ra mình đã bỏ lỡ cơ hội trúng số: ``\textbf{Khônggggggggggg!!} Vậy ra nếu lúc nãy em không nhảy vào phá Matsuri-chan, thì người đi đã là em rồi sao!?''

Cậu Tổng gật đầu, khẽ thở dài, giọng như một người đang an ủi một đứa trẻ vừa làm vỡ đồ chơi: ``Ít nhất thì Mat-chan cũng thoát khỏi đây rồi nhỉ\dots\ Còn chúng ta, chúng ta vẫn còn ở đây với đống chông sắt và con sư tử.''

Căn phòng trở nên tĩnh lặng hơn bao giờ hết. Chỉ còn lại Haato đang quỳ gục xuống, ôm mặt lẩm bẩm gì đó không rõ, trong khi con sư tử vẫn đứng đó, thỉnh thoảng gầm nhẹ, còn Kuon chỉ biết ngồi khoanh tay nhìn cô từ phía sau, đôi mắt đầy sự cảm thông nhưng cũng không kém phần bất lực.

\img{5-10g}

Một lát sau, cô nàng tóc vàng ngẩng đầu lên, mặt đỏ bừng vì tức giận, nắm đấm siết chặt: ``Tất cả là lỗi của Matsuri-chan! Nếu bà ấy không nhảy vào, thì em đã được đi rồi!''

Kuon gõ nhẹ vào trán cô, giọng như một người anh đang nhắc nhở em gái: ``Lỗi của em thì có. Em cũng làm y hệt với em ấy trước đó còn gì. Em bịt miệng cô ấy và gọi sư tử trước. Ăn miếng trả miếng thôi, không có gì phải bất công.''

Haato vùng vẫy, giọng đầy ấm ức: ``Nhưng nếu bà ấy không nhảy vào phá lời em thì đâu ra nông nỗi này! Em đã có thể ước một cách thông minh hơn!''

``Thế thì nếu em không rải chông ngay từ đầu thì đâu có nông nỗi này? Đúng không?''

Haato giật mình, nhưng nhanh chóng đáp lại, giọng vẫn đầy bướng bỉnh: ``Nhưng mà không làm thế thì chán lắm-chama!! Nếu không có chông, thì đâu có vui?''

Nói rồi, cô hậm hực quay mặt đi, hai tay khoanh trước ngực, như một đứa trẻ đang giận dỗi\dots

\dots\ nhưng chợt khựng lại. Đôi mắt cô mở to khi nhìn thấy một thứ gì đó lấp ló giữa đống chông sắt.

``Ể?''

Ngay giữa đống chông sắc nhọn, có một thứ gì đó trông không giống chông chút nào. Một sinh vạt nhỏ nhắn, đen đúa, tròn vo, và đang nhúc nhích một cách kỳ lạ. Haato khom xuống nhìn kỹ hơn\dots\ rồi mắt sáng rỡ như vừa khám phá ra một điều thú vị.

``Đó là con bồ hóng mà!''

Vừa nghe vậy, Kuon ngẩng đầu lên nhìn thử. Sau một hồi quan sát lâu, cậu cũng gật gù, giọng như một nhà động vật học nghiệp dư: ``Quả nhiên đúng là con bồ hóng thật. Không biết có phải cùng loài mà đã làm Mel-chan thành con mực\footnote{Xem lại {\flaregothic ÉPISODE 15}, quyển số Một.} không nữa\dots\ Nhìn nó trông khá quen.''

Con bồ hóng nhỏ bé ấy quay lại, nhìn cô nàng Trung học một lúc lâu, rồi bất ngờ giơ một tấm bảng lên; tấm bảng màu trắng với dòng chữ đen to tướng, thốt ra câu mà cô không ngờ tới: ``Ừm, chính xác!''

\img{5-10h}

Haato chớp mắt đầy ngạc nhiên, giọng như vừa thấy một con thú biết nói: ``Nó\dots\ nó biết nói à? Và còn biết đọc suy nghĩ của mình?''

Trong khi đó, Kuon chỉ biết lặng lẽ thở dài, đưa tay xoa trán: ``Lại cái trò gì nữa đây\dots\ Một con bồ hóng biết đọc suy nghĩ? Chắc chắn đây là một câu chuyện cổ tích khác.''

Nhưng trước khi cậu kịp nói thêm, con bồ hóng đột nhiên tiến đến gần nàng Song tính, đôi mắt đen láy sáng lên, giọng nói đầy phấn khích như một nhân viên du lịch đang chào mời khách hàng: ``Chúc mừng bạn! Bạn đã trúng một chuyến du lịch đến một hòn đảo ở phía Nam!''

Haato đứng sững, mắt mở to: ``Ể? Thật ư!?'' mắt cô sáng lên như hai ngọn đèn pha, cô la lên: ``\textbf{Hoan hôooooo!}''

Nhưng câu nói chưa kịp hoàn tất, một luồng sáng khác bỗng bao trùm lấy cô. Một vòng xoáy kỳ lạ xuất hiện dưới chân cô, và cô bắt đầu bị hút vào.

``N-Này!? Lần này là gì vậy!?''

Cậu quản lý trợn mắt nhìn theo, bất lực thốt lên, giọng như một người đã từ bỏ việc cố gắng hiểu chuyện: ``Lại nữa hả trời\dots\ Hôm nay là ngày gì mà ai cũng biến mất thế này?''

\begin{center}
  \uline{Một hòn đảo nào đó ở phía Nam.}
\end{center}

\img{5-10i}

Haato run cầm cập, mắt rưng rưng, ôm lấy cánh tay trong tuyệt vọng. Cô co ro trong lớp áo mỏng manh, hơi thở bốc thành làn khói trắng trong không khí lạnh giá: ``Haachama-chama\~{}\dots\ Hức\dots\ Lạnh quá\dots\ Đây là đảo phía Nam nào mà toàn tuyết với băng thế này\dots\ Em tưởng phía Nam là nóng chứ\dots''

Một cơn gió lạnh buốt quét qua, tuyết trắng phủ kín mặt đất, và những tảng băng trôi lơ lửng trên mặt nước xám xịt. Cô nhìn quanh, chỉ thấy toàn băng tuyết và một vài chú chim cánh cụt đang đứng nhìn cô với vẻ mặt tò mò.

Bên cạnh cô, một con chim cánh cụt nhỏ đang ngước nhìn cô, tỏ vẻ vui mừng (?) vì cuối cùng đã có một vị khách. Nó vẫy vẫy cánh như thể đang chào đón.

Ngay lúc đó, một cái bóng lớn bay lượn trên bầu trời. Một con chim to lớn nào đó sà xuống, đáp ngay trước mặt Haato. Nó cất giọng đầy hào hứng, như một người dẫn chương trình du lịch: ``Chào mừng đến với Nam Cực! Nơi đây có nhiệt độ trung bình -50 độ C, và chúng tôi có những chuyến tham quan băng tuyết cực kỳ thú vị!''
%==========================================TRUYỆN PHỤ=========================================
\omk

\begin{center}
  \uline{Văn phòng}
\end{center}

Cậu Tổng vẫn đang sững sờ sau cú ``dịch chuyển'' của con bồ hóng. Nó cũng đã biến mất ngay sau khi Haato bị cuốn đi. Giờ đây, trong phòng chỉ còn lại cậu, và những chông sắt rải đầy khắp văn phòng, trông chẳng khác nào một bãi chiến trường sau trận đánh.

``\dots\ Giờ có mỗi mình thôi nhỉ. Không biết cây đèn thần này có thật sự ứng nghiệm không nữa?''

Cậu tò mò cầm lấy cây đèn thần ở trên bàn và xoa nhẹ lên phần hông. Lập tức, một luồng sáng lóe lên, nhưng người xuất hiện không phải là Suisei giống như hai lần trước\dots

``Himemori-na\~{}!''

\dots\ mà lại là Luna. Cậu tròn xoe mắt ngạc nhiên, suýt làm rơi cây đèn: ``Luna-chan!? Anh tưởng là Sui-chan sẽ ra chứ?''

Luna liếc xuống ống vòi đèn, giọng đầy vẻ thông thái của một người vừa phát hiện ra bí mật: ``Chị ấy mợt quá nên ngủ mất tiu rồi-nanora. Sau khi nhờ xong hai điều ước, chị ấy kịt xức và ngủ gục luôn trong đèn rồi.''

Cậu nheo mắt nhìn bên trong, thấy nàng Ca sĩ đã thiếp đi từ khi nào, đầu gục xuống vai, tay vẫn còn cầm tấm bảng Akinator. Cậu thở dài: ``Vậy à\dots\ Em là cũng thần đèn Akinator đó sao?''

Nàng Công chúa nghiêng đầu khó hiểu, đôi mắt long lanh như một đứa trẻ vừa nghe một từ mới: ``A-kê-ne-tơ? Đúng là em với Suisei-paisen là thần đèn, nhưng mà nó là cái gì zậy na?''

Thấy rằng có giải thích cho cô cũng sẽ không hiểu, cậu thở dài một tiếng: ``\dots\ Thôi kệ. Vậy em có thể giúp anh chứ?''

Cô hí hửng đáp, giơ hai tay lên như một vị thần đèn thực thụ: ``Tất nhin rùi! Em sẽ cho anh ba điều ức!''

Kuon nhướng mày ngạc nhiên với lời đề nghị của nàng ``thần đèn'' này: ``Ủa? Không phải là một à?''

Luna cũng nhìn lại cậu với con mắt rất đỗi ngạc nhiên, như thể cậu vừa hỏi một câu vô cùng ngớ ngẩn: ``Hể? Em tửng là câu chuỵn về cây đèn thừn là sẽ theo như zậy chớ-nora?''

Thấy rằng lời đề nghị của cô nàng hào phóng hơn hẳn so với Suisei, cậu cũng chỉ chẹp miệng: ``Cứ cho là vậy.'' Rồi cậu thắc mắc: ``Anh cũng hơi ngạc nhiên là em sẽ là người ban điều ước đấy. Bình thường Công chúa Luna-chan mới là người ước mà?''

Luna phụng phịu nhìn cậu, hai má phồng lên như một đứa trẻ đang giận dỗi: ``Hôm nay em mún đổi vai khum được sao? Đúng là em trẻ con thựt đấy, nhưng mà thi thoản cũng phải thay đủi một tí chứ na?''

Kuon cảm thán, giọng đầy vẻ bất lực: ``\dots\ Vãi đạn. Một cô công chúa 0 tuổi lại muốn làm thần đèn. Mọi thứ đảo lộn tùng phèo hết rồi.'' Cậu xoa cằm, suy nghĩ một lúc rồi nói: ``Điều đầu tiên là: hãy trả lại văn phòng như cũ.''

Nàng Kẹo mỉm cười, giơ tay lên, giọng đầy tự tin: ``Được thôi na\~{}!''

Chỉ trong chớp mắt, toàn bộ những chiếc chông sắt biến thành mây, hai cái ghế sô-pha được trượt về vị trí cũ, căn phòng nhanh chóng trở lại trạng thái nguyên bản. Hai dải cầu vồng thoáng qua trong không trung trước khi dần tan biến, để lại một văn phòng sạch sẽ như chưa từng có chuyện gì xảy ra.

Cậu Tổng nhìn quanh, rồi ngạc nhiên hỏi: ``Nhanh vậy à?''

Luna mỉm cười, tay đặt lên eo, giọng như một người đã làm việc này cả đời: ``Mất một lúc là se bín mớt thui na\~{}. Phép thụt của em lúc nào cũng nhanh, trừ khi em bị đói á\~{}''

Cô vỗ tay một cái rồi tiếp tục: ``Điều ức thứ hai của anh là gì?''

Kuon suy nghĩ trong chốc lát, rồi đáp: ``Đưa anh về nhà đi.''

Luna giơ tay lên, hô một tiếng ``na\~{}'' đầy vui vẻ\dots

\begin{center}
  \uline{Nhà Hikawa.}
\end{center}

\dots\ và ngay lập tức, khung cảnh xung quanh cậu thay đổi hoàn toàn.

Không còn văn phòng. Không còn chông sắt hay ghế sô-pha. Chỉ còn lại căn phòng nhỏ quen thuộc của cậu; một căn phòng trên tầng năm ở khu sầm uất của quận Kyuukami, với những bức tường màu kem và chiếc giường đơn quen thuộc. Ánh đèn vàng ấm áp từ chiếc đèn bàn chiếu xuống, và từ bên ngoài, tiếng xe cộ vọng vào như một bản nhạc đô thị.

Từ trong bếp, một giọng nói quen thuộc cất lên, đầy vẻ ngạc nhiên: ``\vie{Hôm nay con về sớm thế? Mới có năm giờ chiều mà.}''

Cậu liếc xuống chiếc đồng hồ đeo tay, thấy rằng nó đã báo năm giờ chiều. Cậu quay đầu lại, thấy bố cậu đang nhìn ra từ phía bếp với chiếc tạp dề trắng trên người, cậu cười trừ đáp: ``\vie{Sớm sủa gì nữa bố. Năm giờ rồi mà. Với cả hôm nay con mệt nên là về sớm ạ!}''

Mẹ cậu cũng từ trong bếp nói vọng ra, giọng đầy yêu thương: ``\vie{Hôm nay cả nhà ăn mì ramen làm lễ Vu lan đấy con, con lên nhà tắm rửa sạch sẽ rồi xuống ăn nha! Có thịt cháy cạnh và trứng luộc nữa đấy.}''

Kuon chỉ đáp lại một tiếng thật dài, giọng đầy hạnh phúc: ``\vie{Vâng\dots}''

Rồi cậu quay sang Luna, người vẫn đang đứng cạnh cậu với vẻ mặt đầy tò mò khi nhìn căn nhà mới: ``Vậy, điều ức cuối cùm của anh là gì na\~{}?''

Không chút do dự, cậu đáp ngay: ``Cùng với hai em ăn cơm tối một bữa.''

Luna mở to mắt, dường như không tin vào điều mình vừa nghe thấy. Cô nhìn cậu với vẻ ngạc nhiên, như thể cậu vừa yêu cầu một điều gì đó quá đỗi bình thường: ``Ể? Anh hơm mún điều ức gì lớn hơn hả?''

Cậu chỉ bật cười, rồi giải thích thêm, giọng ấm áp: ``Ăn càng đông càng vui mà, đúng chứ? Với cả, coi như hôm nay nhà anh khao mừng lễ Vu lan. Có thêm người thì bữa cơm càng thêm vui.''

Luna chớp mắt một lúc, như thể đang suy nghĩ xem đây có phải là một điều ước ``đáng giá'' hay không. Nhưng rồi cô mỉm cười, gật đầu thật mạnh, vầng hào quang trên đầu xoay nhẹ: ``Đực thôi, nanora\~{}!''

Ngay lập tức, một luồng sáng xuất hiện bên cạnh cô, và một bóng hình quen thuộc dần dần hiện ra: Suisei, người vẫn còn đang ngủ say, tay vẫn ôm tấm bảng Akinator.

Kuon đành phải đưa Suisei lên đệm, cố gắng không làm đánh thức cô dậy. Anh nhẹ nhàng đặt cô xuống, kéo chiếc chăn mỏng đắp lên người cô, rồi quay sang Luna: ``Em ở đây trông Sui-chan, để anh lên nhà nói với bố mẹ.''

Cậu chạy lên tầng bếp, nơi mà bố mẹ cậu vẫn đang tất bật chuẩn bị bữa tối. Mùi thơm của nước dùng ramen lan tỏa khắp căn bếp, và những tiếng dao thái rau củ vang lên đều đều.

``\vie{Bố mẹ ơi, làm cho bọn con thêm hai suất nữa!}''

Ryujin và Kaien cùng quay lại, vẻ mặt ngạc nhiên. Mẹ cậu hỏi, tay vẫn đang cầm chiếc muôi: ``\vie{Hai suất? Sao vậy con? Không phải hôm nay chỉ có nhà mình à?}''

``\vie{Bạn con đến ăn chung ạ,}'' cậu đáp, giọng đầy háo hức, ``\vie{hai người bạn đồng nghiệp của con ạ. Một người là thần đèn, một người là công chúa.}''

Bố cậu bật cười, giọng đầy vẻ hài hước: ``\vie{Thế thì tốt! Hôm nay có khách quý, phải đãi món thật ngon mới được. Để bố làm thêm ít chả cá và trứng rán nữa.}''

Mẹ cậu cũng vui vẻ gật đầu, lau tay vào tạp dề: ``Được rồi, con lên gọi bạn xuống với chuẩn bị chỗ ngồi đi, sắp dọn cơm rồi. Nhớ lấy thêm bát đũa nhé!''

Cậu gật đầu, quay trở lại phòng khách. Luna đang ngồi vắt vẻo trên ghế, trong khi Suisei vừa mới tỉnh dậy, mắt vẫn còn mơ màng như một con mèo vừa ngủ dậy.

Sau khi Suisei vừa mới tỉnh dậy, mắt vẫn còn mơ màng, cô nàng ngáp dài, không quên vươn vai một cái thật to. Kuon ngồi cạnh, nhìn cô với ánh mắt cười mỉm, như thể biết chắc cô ấy sẽ không tỉnh táo cho lắm. Suisei dụi mắt, đưa tay rót thêm một ít trà, rồi khẽ nhấp một ngụm.

``Bộ mặt em có gì đó sao?'' nàng Sao chổi nhíu mày, nhưng giọng nói vẫn lười biếng và không thiếu phần quyến rũ.

``Đâu có,'' Kuon cười đáp, ``chỉ là\dots\ anh thấy em dễ thương quá thôi. Ngái ngủ trông còn dễ thương hơn cả lúc thức.''

Suisei lườm cậu một cái, nhưng nụ cười vẫn không thể giấu được. Cô cố gắng tỏ ra giận dữ, nhưng đôi mắt long lanh đã phản bội cô: ``Anh đừng có tưởng là nói thế sẽ làm em nguôi giận nhé.''

Cả hai im lặng một lúc, thưởng thức bầu không khí yên tĩnh của căn nhà, với ánh đèn ấm áp đang chiếu xuống, cùng hương trà nhẹ nhàng lan tỏa trong không khí.

\ct

Buổi tối hôm đó, trong căn nhà ấm cúng, gia đình cậu cùng với hai thần đèn của cây đèn thần đã có một bữa mì ramen ngon lành. Tiếng cười nói rôm rả, bầu không khí vui vẻ lan tỏa khắp phòng ăn.

Một chút gió se lạnh từ bên ngoài thổi qua cửa sổ, nhưng không gian trong phòng vẫn ấm áp nhờ sự hiện diện của bữa tối đang bày ra. Mùi mì ramen thơm lừng, cùng với sự ấm áp của gia đình, làm cho cậu cảm thấy một cảm giác bình yên hiếm có.

Giữa lúc đang thưởng thức tô mì nóng hổi, cậu Tổng chợt dừng lại, gõ đũa vào miệng suy nghĩ: ``Hình như anh cảm thấy anh đang quên điều gì đó\dots''

Suisei vẫn còn hơi ngái ngủ, nhấp một ngụm nước rồi hờ hững đáp lại như chẳng hề liên quan: ``Nếu anh quên nó thì chắc không quan trọng lắm đâu. Nếu quan trọng thì anh đã nhớ rồi.''

Cậu chớp mắt vài lần, rồi nhún vai: ``\dots Ừ, cũng có lý. Chắc là đầu anh quá tải rồi\dots''

Và thế là, họ tiếp tục thưởng thức món ăn.

Trong khi đó, Luna thì reo lên với giọng đầy phấn khích, hai tay cầm bát mì: ``Bát mì ramen ngon thựt đấy na\~{}!! Kuoncha-senpai, anh muốn thử bát của em hơm? Có thịt cháy cạnh với trứng lục đấy!''

Cậu nhìn cô, rồi nhìn bát mì, và mỉm cười: ``Cảm ơn em, nhưng anh có bát của mình rồi. Cứ ăn đi, kẻo nguội.''

\begin{center}
  \uline{Nam Cực.}
\end{center}

Giữa cơn bão tuyết lạnh cắt da cắt thịt, một bóng dáng nhỏ bé lầm lũi bước đi, cố gắng tìm kiếm sự trợ giúp. Tuyết rơi dày đặc, gió rít lên từng cơn, và nhiệt độ đã xuống dưới -40 độ C.

``Cứu tôi với\dots\ Có ai ở đây không\dots?''

Giọng Haato vang lên yếu ớt giữa không gian hoang vu. Gió lạnh táp vào mặt, khiến cô rùng mình liên tục, đôi tay đã tê cóng vì lạnh.

``Tại sao mình lại bị như vậy chứ\dots''

Vừa dứt lời, cô hắt hơi một cái rõ to, tiếng hắt hơi vang xa trong không gian tĩnh mịch.

Về sau, may mắn thay, một đoàn cứu hộ đã tìm thấy cô và đưa cô lên trực thăng trở về nước an toàn. Khi được hỏi tại sao cô lại ở Nam Cực, cô chỉ biết khóc và nói: ``Tôi chỉ muốn rời khỏi đây thôi mà\dots''

\begin{center}
  \uline{Một quán mì ramen ở Kawachi.}
\end{center}

Cùng thời điểm mà gia đình Hikawa đang dùng bữa mì ramen, ở một nơi khác tại chính thành phố mà gia đình cậu ở, cũng có một người khác đang thưởng thức món mì ramen ngon lành.

Mùi nước dùng thơm lừng lan tỏa khắp quán ăn. Các thực khách đang dùng bữa thì đột nhiên chú ý đến một cô gái tóc nâu đang ngồi ở góc quán, trước mặt là một đống bát mì đã ăn sạch.

Matsuri hớn hở giơ tay gọi thêm một bát nữa, giọng như một đứa trẻ vừa tìm thấy kho báu: ``Hmmm\~{} Ngon quá! Cho tôi thêm một bát nữa nhé! Lần này thêm trứng luộc và chả cá!''

Trước ánh mắt ngỡ ngàng của những người xung quanh, cô đã đánh chén xong bốn bát mì cỡ lớn một cách đầy ngoạn mục. Chủ quán nhìn cô với vẻ mặt vừa kinh ngạc vừa ngưỡng mộ.

Ở một góc khác của quán, hai người bạn cũ của Kuon cũng đang ngồi theo dõi màn trình diễn ẩm thực có một không hai này.

Yuishi ngạc nhiên thốt lên, giọng như không tin nổi vào mắt mình: ``\vie{Quắt đờ phắc\dots\ Nó mới ăn hết bốn bát rồi mà vẫn gọi thêm ư!?}''

Sekitei nhìn chằm chằm vào đống bát trước mặt Matsuri, lẩm bẩm: ``\vie{Nó ăn đến bát thứ năm rồi cơ á\dots? Mà nó vẫn gọi thêm. Chắc nó có dạ dày không đáy luôn rồi\dots}''

Cuối cùng, sau khi đánh chén hết bảy bát mì, nàng Lễ hội vỗ bụng thoả mãn, vui vẻ trả tiền rồi bắt xe về văn phòng công ty với một cái bụng no căng, quên hết sự đời. Trên xe, cô còn lẩm bẩm: ``Ngày mai lại quay lại quán này\dots\ Món ramen ở đây ngon thật đấy.''

\il

Ngày lễ Vu lan năm nay đúng là kỳ lạ thật! Một ngày bắt đầu với đầy chông sắt, kết thúc với những bát mì ramen và những cuộc phiêu lưu không tưởng. Và có lẽ, đó là cách mà \hll\ ăn mừng một ngày lễ truyền thống theo cách riêng của họ.

\end{document}